\documentclass[11pt, ukrainian]{article}
\setlength\textwidth{180mm} \setlength\textheight{277mm}
\setlength\voffset{-38mm} \setlength\hoffset{-28mm}
%\setlength\parindent{0mm}
\pagestyle{empty}
\usepackage[T2A]{fontenc}
\usepackage[utf8]{inputenc}
\usepackage[ukrainian]{babel}
\begin{document}
{\Large \center  Варіанти завдань до лабораторної роботи №4\par 2020/21 н.р.\par \bigskip\bigskip}

Порядок полів у записах та інформацію додаткового файлу фіксовано у варіанті.

Для порівняння числових значень використовується традиційний порядок.

Для порівняння рядкових значень використовується лексикографічний порядок. 

У випадку сортування за кількома полями результуюче сортування будується 
за порядками окремих полів за допомогою конструкції лексикографічного порядку. 
Послідовність полів, за якими має відбуватися сортування, задається в умові. 
Якщо не вказано інше, то поле сортується за зростанням.

Якщо умова вимагає знайти $k$ най… і $k$-е місце займають кілька, то слід виводити всіх їх.

\newpage
\par \bigskip

{\center Варіант  \No \textbf { 1 }\par}
\bigskip
%type = absense
%variant = 22
Обробляються дані семестрового журналу перевірок відвідування занять студентами деканатом (фіксуються відсутні). 

\underline{Записи основного файлу містять поля}: аудиторія, вид занять, прізвище, ім’я, предмет, навчальний тиждень, день тижня, пара, курс, код групи.

\underline{У допоміжному файлі наявні ключі}: кількість лекцій, кількість записів.

\underline{Знайти} всіх студентів, що пропустили пар не більше ніж в середньому. Вивести по кожному з них інформацію:\par
{\noindent}--- на першому рядку:\par
прізвище, ім’я, кількість пропущених пар, середню кількість пропусків по всіх студентах (округлена до 1 знаку);

{\noindent}--- на наступних рядках, починаючи з табуляції, вивести пропуски студента (по одному на рядок):\par
тиждень, день, пара, аудиторія, предмет, вид занять

{\noindent}у такому сортуванні: вид занять, предмет, тиждень, день, пара.

%Студентів сортувати: кількість пропусків, прізвище, ім’я.

\bigskip\textit{Обмеження}

Навчальний тиждень може приймати цілі значення від 1 до 19.
День тижня є цілим від 1 до 5. Пара є цілим від 1 до 4.
Курс є цілим від 1 до 4. Аудиторія є додатнім цілим.
Вид занять може приймати значення:
lect --- лекція, P --- практичне, S --- семінар, lab --- лабораторне.

Решта полів є непорожніми рядками.
Рядки починатися та закінчуватися білими символами не можуть.

Групи кодуються в межах курсу.
Код групи складається не більш ніж з 5 символів.
Може містити літери алфавіту, десяткові цифри та дефіс.

Предмет є рядком, може мати від 4 до 35 символів включно.
Може містити літери алфавіту, десяткові цифри, апостроф, дефіс, пробіл.
Предмети різних курсів вважати різними навіть за однакових назв.

Прізвище, ім’я та по-батькові (за наявності у варіанті)
можуть мати від 4 до 30, 28, 30 відповідно символів включно кожне.
Можуть містити літери алфавіту, апостроф, пробіл та дефіс.
Студенти однозначно ідентифікуються прізвищем, ім’ям та по-батькові (за наявності у варіанті).
Жоден студент не вчиться одночасно на різних курсах чи в різних групах.

Види занять сортуються так: лекція < семінар < практичне < лабораторне.

%\bigskip%\textit{Для деяких варіантів може знадобитись}


\par
\newpage
{\center Варіант  \No \textbf { 2 }\par}
\bigskip
%type = absense
%variant = 30
Обробляються дані семестрового журналу перевірок відвідування занять студентами деканатом (фіксуються відсутні). 

\underline{Записи основного файлу містять поля}: навчальний тиждень, ім’я, предмет, вид занять, по-батькові, прізвище, день тижня, пара, курс, код групи, аудиторія.

\underline{У допоміжному файлі наявні ключі}: найменший номер аудиторії, сума номерів пар.

\underline{Знайти} всіх студентів, що пропустили практичних більше ніж в середньому (практичних було пропущено студентами, що пропускали практичні заняття). Вивести по кожному з них інформацію:\par
{\noindent}--- на першому рядку:\par
прізвище, ім’я по-батькові, кількість пропущених пар, кількість пропущених практичних;

{\noindent}--- на наступних рядках, починаючи з табуляції, вивести пропуски студента (по одному на рядок):\par
предмет, вид занять, тиждень, день, пара

{\noindent}у такому сортуванні: предмет, вид занять, тиждень, день, пара.

%Студентів сортувати: кількість пропусків практичних (за спаданням), кількість пропусків лекцій, прізвище, ім’я, по-батькові.

\bigskip\textit{Обмеження}

Навчальний тиждень може приймати цілі значення від 1 до 17.
День тижня є цілим від 1 до 5. Пара є цілим від 1 до 4.
Курс є цілим від 1 до 4. Аудиторія є додатнім цілим.
Вид занять може приймати значення:
Le --- лекція, p --- практичне, сем. --- семінар, Lab --- лабораторне.

Решта полів є непорожніми рядками.
Рядки починатися та закінчуватися білими символами не можуть.

Групи кодуються в межах курсу.
Код групи складається не більш ніж з 6 символів.
Може містити літери алфавіту, десяткові цифри та дефіс.

Предмет є рядком, може мати від 3 до 25 символів включно.
Може містити літери алфавіту, десяткові цифри, апостроф, дефіс, пробіл.
Предмети різних курсів вважати різними навіть за однакових назв.

Прізвище, ім’я та по-батькові (за наявності у варіанті)
можуть мати від 3 до 21, 29, 23 відповідно символів включно кожне.
Можуть містити літери алфавіту, апостроф, пробіл та дефіс.
Студенти однозначно ідентифікуються прізвищем, ім’ям та по-батькові (за наявності у варіанті).
Жоден студент не вчиться одночасно на різних курсах чи в різних групах.

Види занять сортуються так: практичне < семінар < лекція < лабораторне.

%\bigskip%\textit{Для деяких варіантів може знадобитись}


\par
\newpage
{\center Варіант  \No \textbf { 3 }\par}
\bigskip
%type = absense
%variant = 49
Обробляються дані семестрового журналу перевірок відвідування занять студентами деканатом (фіксуються відсутні). 

\underline{Записи основного файлу містять поля}: предмет, день тижня, пара, курс, прізвище, ім’я, аудиторія, вид занять, навчальний тиждень.

\underline{У допоміжному файлі наявні ключі}: кількість пропусків третіх пар, кількість різних предметів.

\underline{Знайти} всі предмети, які пропускалися тільки по понеділках. Вивести по кожному з них інформацію:\par
{\noindent}--- на першому рядку:\par
курс, предмет, кількість пропусків;

{\noindent}--- на наступних рядках, починаючи з табуляції, вивести пропуски предмета (по одному на рядок):\par
прізвище, ім’я, тиждень, кількість пропусків

{\noindent}у такому сортуванні: прізвище, ім’я, тиждень.

%Предмети сортувати: предмет, курс.

\bigskip\textit{Обмеження}

Навчальний тиждень може приймати цілі значення від 1 до 20.
День тижня є цілим від 1 до 5. Пара є цілим від 1 до 4.
Курс є цілим від 1 до 4. Аудиторія є додатнім цілим.
Вид занять може приймати значення:
Лекц. --- лекція, п --- практичне, Сем. --- семінар, лаб. --- лабораторне.

Решта полів є непорожніми рядками.
Рядки починатися та закінчуватися білими символами не можуть.

Групи кодуються в межах курсу.
Код групи складається не більш ніж з 4 символів.
Може містити літери алфавіту, десяткові цифри та дефіс.

Предмет є рядком, може мати від 5 до 27 символів включно.
Може містити літери алфавіту, десяткові цифри, апостроф, дефіс, пробіл.
Предмети різних курсів вважати різними навіть за однакових назв.

Прізвище, ім’я та по-батькові (за наявності у варіанті)
можуть мати від 5 до 26, 23, 27 відповідно символів включно кожне.
Можуть містити літери алфавіту, апостроф, пробіл та дефіс.
Студенти однозначно ідентифікуються прізвищем, ім’ям та по-батькові (за наявності у варіанті).
Жоден студент не вчиться одночасно на різних курсах чи в різних групах.

Види занять сортуються так: практичне < лекція < семінар < лабораторне.

%\bigskip%\textit{Для деяких варіантів може знадобитись}


\par
\newpage
{\center Варіант  \No \textbf { 4 }\par}
\bigskip
%type = absense
%variant = 27
Обробляються дані семестрового журналу перевірок відвідування занять студентами деканатом (фіксуються відсутні). 

\underline{Записи основного файлу містять поля}: предмет, курс, прізвище, ім’я, вид занять, навчальний тиждень, день тижня, пара, аудиторія.

\underline{У допоміжному файлі наявні ключі}: кількість пропусків лекцій, загальна кількість записів.

\underline{Знайти} всіх студентів, що пропускали лабораторні заняття. Вивести по кожному з них інформацію:\par
{\noindent}--- на першому рядку:\par
кількість пропущених пар, кількість пропущених лабораторних, відсоток пропусків лабораторних серед всіх пропусків студента (округлений до цілого), прізвище, ім’я,;

{\noindent}--- на наступних рядках, починаючи з табуляції, вивести пропуски лабораторних занять студентом (по одному на рядок):\par
предмет, тиждень, день, пара, аудиторія

{\noindent}у такому сортуванні: предмет, тиждень, день, пара.

%Студентів сортувати: прізвище, ім’я.

\bigskip\textit{Обмеження}

Навчальний тиждень може приймати цілі значення від 1 до 16.
День тижня є цілим від 1 до 5. Пара є цілим від 1 до 4.
Курс є цілим від 1 до 4. Аудиторія є додатнім цілим.
Вид занять може приймати значення:
Л --- лекція, Практ. --- практичне, s --- семінар, Лаб. --- лабораторне.

Решта полів є непорожніми рядками.
Рядки починатися та закінчуватися білими символами не можуть.

Групи кодуються в межах курсу.
Код групи складається не більш ніж з 4 символів.
Може містити літери алфавіту, десяткові цифри та дефіс.

Предмет є рядком, може мати від 6 до 32 символів включно.
Може містити літери алфавіту, десяткові цифри, апостроф, дефіс, пробіл.
Предмети різних курсів вважати різними навіть за однакових назв.

Прізвище, ім’я та по-батькові (за наявності у варіанті)
можуть мати від 6 до 29, 30, 24 відповідно символів включно кожне.
Можуть містити літери алфавіту, апостроф, пробіл та дефіс.
Студенти однозначно ідентифікуються прізвищем, ім’ям та по-батькові (за наявності у варіанті).
Жоден студент не вчиться одночасно на різних курсах чи в різних групах.

Види занять сортуються так: лабораторне < практичне < лекція < семінар.

%\bigskip%\textit{Для деяких варіантів може знадобитись}


\par
\newpage
{\center Варіант  \No \textbf { 5 }\par}
\bigskip
%type = absense
%variant = 38
Обробляються дані семестрового журналу перевірок відвідування занять студентами деканатом (фіксуються відсутні). 

\underline{Записи основного файлу містять поля}: аудиторія, день тижня, пара, курс, вид занять, навчальний тиждень, предмет, прізвище, ім’я.

\underline{У допоміжному файлі наявні ключі}: кількість різних предметів, довжина найкоротшого прізвища.

\underline{Знайти} всі предмети, що пропускалися найбільше. Вивести по кожному з них інформацію:\par
{\noindent}--- на першому рядку:\par
курс, предмет, кількість пропусків, кількість пропусків лекцій;

{\noindent}--- на наступних рядках, починаючи з табуляції, вивести пропуски предмета (по одному на рядок):\par
прізвище, ім’я, тиждень, день, пара, вид занять

{\noindent}у такому сортуванні: вид занять, прізвище, ім’я, тиждень, день, пара.

%Предмети сортувати: кількість пропусків(за спаданням), курс, предмет.

\bigskip\textit{Обмеження}

Навчальний тиждень може приймати цілі значення від 1 до 18.
День тижня є цілим від 1 до 5. Пара є цілим від 1 до 4.
Курс є цілим від 1 до 4. Аудиторія є додатнім цілим.
Вид занять може приймати значення:
Lecture --- лекція, ПР --- практичне, sem --- семінар, ЛАБ --- лабораторне.

Решта полів є непорожніми рядками.
Рядки починатися та закінчуватися білими символами не можуть.

Групи кодуються в межах курсу.
Код групи складається не більш ніж з 5 символів.
Може містити літери алфавіту, десяткові цифри та дефіс.

Предмет є рядком, може мати від 4 до 21 символів включно.
Може містити літери алфавіту, десяткові цифри, апостроф, дефіс, пробіл.
Предмети різних курсів вважати різними навіть за однакових назв.

Прізвище, ім’я та по-батькові (за наявності у варіанті)
можуть мати від 4 до 24, 21, 26 відповідно символів включно кожне.
Можуть містити літери алфавіту, апостроф, пробіл та дефіс.
Студенти однозначно ідентифікуються прізвищем, ім’ям та по-батькові (за наявності у варіанті).
Жоден студент не вчиться одночасно на різних курсах чи в різних групах.

Види занять сортуються так: лабораторне < практичне < лекція < семінар.

%\bigskip%\textit{Для деяких варіантів може знадобитись}


\par
\newpage
{\center Варіант  \No \textbf { 6 }\par}
\bigskip
%type = absense
%variant = 28
Обробляються дані семестрового журналу перевірок відвідування занять студентами деканатом (фіксуються відсутні). 

\underline{Записи основного файлу містять поля}: ім’я, по-батькові, прізвище, курс, код групи, аудиторія, вид занять, день тижня, пара, навчальний тиждень, предмет.

\underline{У допоміжному файлі наявні ключі}: кількість студентів, найбільша довжина прізвища.

\underline{Знайти} всіх студентів, що пропускали лекції, але не пропускали практичні та лабораторні заняття. Вивести по кожному з них інформацію:\par
{\noindent}--- на першому рядку:\par
кількість пропусків, прізвище, ім’я, по-батькові;

{\noindent}--- на наступних рядках, починаючи з табуляції, вивести пропуски студента (по одному на рядок):\par
предмет, тиждень, кількість пропусків лекцій, кількість пропусків практичних, кількість пропусків лабораторних, кількість пропусків семінарів, загальна кількість пропусків

{\noindent}у такому сортуванні: предмет, тиждень.

%Студентів сортувати: кількість пропусків(за спаданням), прізвище, ім’я, по-батькові.

\bigskip\textit{Обмеження}

Навчальний тиждень може приймати цілі значення від 1 до 14.
День тижня є цілим від 1 до 5. Пара є цілим від 1 до 4.
Курс є цілим від 1 до 4. Аудиторія є додатнім цілим.
Вид занять може приймати значення:
L --- лекція, практ. --- практичне, Sem --- семінар, lab --- лабораторне.

Решта полів є непорожніми рядками.
Рядки починатися та закінчуватися білими символами не можуть.

Групи кодуються в межах курсу.
Код групи складається не більш ніж з 6 символів.
Може містити літери алфавіту, десяткові цифри та дефіс.

Предмет є рядком, може мати від 6 до 39 символів включно.
Може містити літери алфавіту, десяткові цифри, апостроф, дефіс, пробіл.
Предмети різних курсів вважати різними навіть за однакових назв.

Прізвище, ім’я та по-батькові (за наявності у варіанті)
можуть мати від 6 до 20, 25, 28 відповідно символів включно кожне.
Можуть містити літери алфавіту, апостроф, пробіл та дефіс.
Студенти однозначно ідентифікуються прізвищем, ім’ям та по-батькові (за наявності у варіанті).
Жоден студент не вчиться одночасно на різних курсах чи в різних групах.

Види занять сортуються так: практичне < лабораторне < лекція < семінар.

%\bigskip%\textit{Для деяких варіантів може знадобитись}


\par
\newpage
{\center Варіант  \No \textbf { 7 }\par}
\bigskip
%type = absense
%variant = 51
Обробляються дані семестрового журналу перевірок відвідування занять студентами деканатом (фіксуються відсутні). 

\underline{Записи основного файлу містять поля}: день тижня, пара, навчальний тиждень, ім’я, вид занять, прізвище, аудиторія, предмет, курс.

\underline{У допоміжному файлі наявні ключі}: загальна кількість предметів, кількість пропусків перших пар.

\underline{Знайти} всі предмети, що по середах пропускалися більше ніж у середньому. Вивести по кожному з них інформацію:\par
{\noindent}--- на першому рядку:\par
курс, предмет, кількість пропусків, кількість пропусків по середах;

{\noindent}--- на наступних рядках, починаючи з табуляції, вивести пропуски предмета (по одному на рядок):\par
прізвище, ім’я, день, вид занять, кількість пропусків

{\noindent}у такому сортуванні: день, вид занять, прізвище, ім’я.

%Предмети сортувати: кількість пропусків по середах, курс, предмет.

\bigskip\textit{Обмеження}

Навчальний тиждень може приймати цілі значення від 1 до 15.
День тижня є цілим від 1 до 5. Пара є цілим від 1 до 4.
Курс є цілим від 1 до 4. Аудиторія є додатнім цілим.
Вид занять може приймати значення:
лекц. --- лекція, pr --- практичне, с. --- семінар, лаб. --- лабораторне.

Решта полів є непорожніми рядками.
Рядки починатися та закінчуватися білими символами не можуть.

Групи кодуються в межах курсу.
Код групи складається не більш ніж з 6 символів.
Може містити літери алфавіту, десяткові цифри та дефіс.

Предмет є рядком, може мати від 3 до 23 символів включно.
Може містити літери алфавіту, десяткові цифри, апостроф, дефіс, пробіл.
Предмети різних курсів вважати різними навіть за однакових назв.

Прізвище, ім’я та по-батькові (за наявності у варіанті)
можуть мати від 3 до 25, 24, 22 відповідно символів включно кожне.
Можуть містити літери алфавіту, апостроф, пробіл та дефіс.
Студенти однозначно ідентифікуються прізвищем, ім’ям та по-батькові (за наявності у варіанті).
Жоден студент не вчиться одночасно на різних курсах чи в різних групах.

Види занять сортуються так: лекція < практичне < семінар < лабораторне.

%\bigskip%\textit{Для деяких варіантів може знадобитись}


\par
\newpage
{\center Варіант  \No \textbf { 8 }\par}
\bigskip
%type = absense
%variant = 32
Обробляються дані семестрового журналу перевірок відвідування занять студентами деканатом (фіксуються відсутні). 

\underline{Записи основного файлу містять поля}: ім’я, день тижня, пара, прізвище, аудиторія, вид занять, курс, предмет, навчальний тиждень.

\underline{У допоміжному файлі наявні ключі}: загальна кількість\ студентів, кількість пропусків\ по п’ятни\-цях.

\underline{Знайти} всіх студентів, що пропускали заняття на перших парах. Вивести по кожному з них інформацію:\par
{\noindent}--- на першому рядку:\par
прізвище, ім’я, кількість пропущених пар, кількість пропущених перших пар;

{\noindent}--- на наступних рядках, починаючи з табуляції, вивести пропуски студента (по одному на рядок):\par
день, пара, предмет, вид занять, кількість пропусків

{\noindent}у такому сортуванні: день, пара, предмет, вид занять.

%Студентів сортувати: кількість пропусків перших пар(за спаданням), кількість пропусків(за спаданням), прізвище, ім’я.

\bigskip\textit{Обмеження}

Навчальний тиждень може приймати цілі значення від 1 до 16.
День тижня є цілим від 1 до 5. Пара є цілим від 1 до 4.
Курс є цілим від 1 до 4. Аудиторія є додатнім цілим.
Вид занять може приймати значення:
Lect --- лекція, Practice --- практичне, Seminar --- семінар, Лаб. --- лабораторне.

Решта полів є непорожніми рядками.
Рядки починатися та закінчуватися білими символами не можуть.

Групи кодуються в межах курсу.
Код групи складається не більш ніж з 4 символів.
Може містити літери алфавіту, десяткові цифри та дефіс.

Предмет є рядком, може мати від 5 до 31 символів включно.
Може містити літери алфавіту, десяткові цифри, апостроф, дефіс, пробіл.
Предмети різних курсів вважати різними навіть за однакових назв.

Прізвище, ім’я та по-батькові (за наявності у варіанті)
можуть мати від 5 до 22, 26, 21 відповідно символів включно кожне.
Можуть містити літери алфавіту, апостроф, пробіл та дефіс.
Студенти однозначно ідентифікуються прізвищем, ім’ям та по-батькові (за наявності у варіанті).
Жоден студент не вчиться одночасно на різних курсах чи в різних групах.

Види занять сортуються так: семінар < лабораторне < практичне < лекція.

%\bigskip%\textit{Для деяких варіантів може знадобитись}


\par
\newpage
{\center Варіант  \No \textbf { 9 }\par}
\bigskip
%type =     absense
%variant = 21
Обробляються дані семестрового журналу перевірок відвідування занять студентами деканатом (фіксуються відсутні). 

\underline{Записи основного файлу містять поля}: ім’я, по-батькові, день тижня, пара, вид занять, аудиторія, прізвище, курс, код групи, навчальний тиждень, предмет.

\underline{У допоміжному файлі наявні ключі}: загальна кількість зафіксованих пропусків, сума курсів.

\underline{Знайти} всіх студентів, які пропустили найбільшу кількість пар. Вивести по кожному з них інформацію:\par
{\noindent}--- на першому рядку:\par
прізвище, ім’я, по-батькові, кількість пропущених пар, кількість пропущених лекцій, кількість пропущених семінарів;

{\noindent}--- на наступних рядках, починаючи з табуляції, вивести пропуски студента (по одному на рядок):\par
тиждень, день, пара, аудиторія, предмет, вид занять

{\noindent}у такому сортуванні: предмет, тиждень, день, пара, вид занять.

%Студентів сортувати: кількість пропусків (за спаданням), кількість пропусків лекцій, прізвище, ім’я, по-батькові.

\bigskip\textit{Обмеження}

Навчальний тиждень може приймати цілі значення від 1 до 14.
День тижня є цілим від 1 до 5. Пара є цілим від 1 до 4.
Курс є цілим від 1 до 4. Аудиторія є додатнім цілим.
Вид занять може приймати значення:
ЛЕ --- лекція, Pr --- практичне, Sem --- семінар, Lab --- лабораторне.

Решта полів є непорожніми рядками.
Рядки починатися та закінчуватися білими символами не можуть.

Групи кодуються в межах курсу.
Код групи складається не більш ніж з 5 символів.
Може містити літери алфавіту, десяткові цифри та дефіс.

Предмет є рядком, може мати від 5 до 22 символів включно.
Може містити літери алфавіту, десяткові цифри, апостроф, дефіс, пробіл.
Предмети різних курсів вважати різними навіть за однакових назв.

Прізвище, ім’я та по-батькові (за наявності у варіанті)
можуть мати від 5 до 27, 20, 25 відповідно символів включно кожне.
Можуть містити літери алфавіту, апостроф, пробіл та дефіс.
Студенти однозначно ідентифікуються прізвищем, ім’ям та по-батькові (за наявності у варіанті).
Жоден студент не вчиться одночасно на різних курсах чи в різних групах.

Види занять сортуються так: лабораторне < лекція < практичне < семінар.

%\bigskip%\textit{Для деяких варіантів може знадобитись}


\par
\newpage
{\center Варіант  \No \textbf { 10 }\par}
\bigskip
%type = absense
%variant = 35
Обробляються дані семестрового журналу перевірок відвідування занять студентами деканатом (фіксуються відсутні). 

\underline{Записи основного файлу містять поля}: курс, код групи, ім’я, прізвище, навчальний тиждень, аудиторія, вид занять, предмет, день тижня, пара.

\underline{У допоміжному файлі наявні ключі}: BOOK, загальна кількість записів.

\underline{Знайти} всіх студентів, що пропускали заняття тільки по п’ятницях. Вивести по кожному з них інформацію:\par
{\noindent}--- на першому рядку:\par
кількість пропусків, прізвище, ім’я;

{\noindent}--- на наступних рядках, починаючи з табуляції, вивести пропуски студента (по одному на рядок):\par
день, пара, загальна кількість пропусків

{\noindent}у такому сортуванні: день, пара.

%Студентів сортувати: кількість пропусків, прізвище, ім’я.

\bigskip\textit{Обмеження}

Навчальний тиждень може приймати цілі значення від 1 до 20.
День тижня є цілим від 1 до 5. Пара є цілим від 1 до 4.
Курс є цілим від 1 до 4. Аудиторія є додатнім цілим.
Вид занять може приймати значення:
Lecture --- лекція, Пр --- практичне, с. --- семінар, ЛАБ --- лабораторне.

Решта полів є непорожніми рядками.
Рядки починатися та закінчуватися білими символами не можуть.

Групи кодуються в межах курсу.
Код групи складається не більш ніж з 4 символів.
Може містити літери алфавіту, десяткові цифри та дефіс.

Предмет є рядком, може мати від 6 до 36 символів включно.
Може містити літери алфавіту, десяткові цифри, апостроф, дефіс, пробіл.
Предмети різних курсів вважати різними навіть за однакових назв.

Прізвище, ім’я та по-батькові (за наявності у варіанті)
можуть мати від 6 до 23, 27, 20 відповідно символів включно кожне.
Можуть містити літери алфавіту, апостроф, пробіл та дефіс.
Студенти однозначно ідентифікуються прізвищем, ім’ям та по-батькові (за наявності у варіанті).
Жоден студент не вчиться одночасно на різних курсах чи в різних групах.

Види занять сортуються так: практичне < семінар < лекція < лабораторне.

%\bigskip%\textit{Для деяких варіантів може знадобитись}


\par
\newpage
{\center Варіант  \No \textbf { 11 }\par}
\bigskip
%type = absense
%variant = 25
Обробляються дані семестрового журналу перевірок відвідування занять студентами деканатом (фіксуються відсутні). 

\underline{Записи основного файлу містять поля}: навчальний тиждень, курс, день тижня, пара, аудиторія, ім’я, вид занять, прізвище, предмет.

\underline{У допоміжному файлі наявні ключі}: суму номерів пар, довжину найкоротшого предмета.

\underline{Знайти} всіх студентів, що пропускали пари тільки перші 3 навчальних тижні. Вивести по кожному з них інформацію:\par
{\noindent}--- на першому рядку:\par
прізвище, ім’я, кількість пропущених пар, кількість пропущених лабораторних;

{\noindent}--- на наступних рядках, починаючи з табуляції, вивести пропуски студента (по одному на рядок):\par
тиждень, день, пара, аудиторія, предмет, вид занять

{\noindent}у такому сортуванні: тиждень, день, пара, предмет, вид занять.

%Студентів сортувати: останній тиждень, на якому були пропуски (за спаданням), прізвище, ім’я.

\bigskip\textit{Обмеження}

Навчальний тиждень може приймати цілі значення від 1 до 18.
День тижня є цілим від 1 до 5. Пара є цілим від 1 до 4.
Курс є цілим від 1 до 4. Аудиторія є додатнім цілим.
Вид занять може приймати значення:
Лекц. --- лекція, п --- практичне, s --- семінар, лаб. --- лабораторне.

Решта полів є непорожніми рядками.
Рядки починатися та закінчуватися білими символами не можуть.

Групи кодуються в межах курсу.
Код групи складається не більш ніж з 5 символів.
Може містити літери алфавіту, десяткові цифри та дефіс.

Предмет є рядком, може мати від 3 до 34 символів включно.
Може містити літери алфавіту, десяткові цифри, апостроф, дефіс, пробіл.
Предмети різних курсів вважати різними навіть за однакових назв.

Прізвище, ім’я та по-батькові (за наявності у варіанті)
можуть мати від 3 до 28, 22, 29 відповідно символів включно кожне.
Можуть містити літери алфавіту, апостроф, пробіл та дефіс.
Студенти однозначно ідентифікуються прізвищем, ім’ям та по-батькові (за наявності у варіанті).
Жоден студент не вчиться одночасно на різних курсах чи в різних групах.

Види занять сортуються так: семінар < лекція < практичне < лабораторне.

%\bigskip%\textit{Для деяких варіантів може знадобитись}


\par
\newpage
{\center Варіант  \No \textbf { 12 }\par}
\bigskip
%type = absense
%variant = 24
Обробляються дані семестрового журналу перевірок відвідування занять студентами деканатом (фіксуються відсутні). 

\underline{Записи основного файлу містять поля}: курс, код групи, навчальний тиждень, предмет, ім’я, прізвище, вид занять, день тижня, пара, аудиторія.

\underline{У допоміжному файлі наявні ключі}: найменша довжина назви предмета, найбільша довжина назви предмета, загальна кількість студентів.

\underline{Знайти} всіх студентів, що пропускали пари тільки з одного предмета. Вивести по кожному з них інформацію:\par
{\noindent}--- на першому рядку:\par
прізвище, ім’я, кількість пропущених пар, кількість пропущених лекцій, кількість пропущених лабораторних, предмет, що пропускався;

{\noindent}--- на наступних рядках, починаючи з табуляції, вивести пропуски студента (по одному на рядок):\par
тиждень, день, пара, аудиторія, вид занять

{\noindent}у такому сортуванні: пара, тиждень, день, вид занять.

%Студентів сортувати: прізвище, ім’я.

\bigskip\textit{Обмеження}

Навчальний тиждень може приймати цілі значення від 1 до 17.
День тижня є цілим від 1 до 5. Пара є цілим від 1 до 4.
Курс є цілим від 1 до 4. Аудиторія є додатнім цілим.
Вид занять може приймати значення:
Л --- лекція, Pr --- практичне, сем. --- семінар, lab --- лабораторне.

Решта полів є непорожніми рядками.
Рядки починатися та закінчуватися білими символами не можуть.

Групи кодуються в межах курсу.
Код групи складається не більш ніж з 6 символів.
Може містити літери алфавіту, десяткові цифри та дефіс.

Предмет є рядком, може мати від 4 до 38 символів включно.
Може містити літери алфавіту, десяткові цифри, апостроф, дефіс, пробіл.
Предмети різних курсів вважати різними навіть за однакових назв.

Прізвище, ім’я та по-батькові (за наявності у варіанті)
можуть мати від 4 до 25, 30, 30 відповідно символів включно кожне.
Можуть містити літери алфавіту, апостроф, пробіл та дефіс.
Студенти однозначно ідентифікуються прізвищем, ім’ям та по-батькові (за наявності у варіанті).
Жоден студент не вчиться одночасно на різних курсах чи в різних групах.

Види занять сортуються так: лекція < лабораторне < практичне < семінар.

%\bigskip%\textit{Для деяких варіантів може знадобитись}


\par
\newpage
{\center Варіант  \No \textbf { 13 }\par}
\bigskip
%type = absense
%variant = 46
Обробляються дані семестрового журналу перевірок відвідування занять студентами деканатом (фіксуються відсутні). 

\underline{Записи основного файлу містять поля}: вид занять, ім’я, предмет, курс, навчальний тиждень, аудиторія, прізвище, день тижня, пара.

\underline{У допоміжному файлі наявні ключі}: кількість пропусків по вівторках, загальна кількість зафіксованих пропусків.

\underline{Знайти} всі предмети, з яких пропускалися тільки лекції. Вивести по кожному з них інформацію:\par
{\noindent}--- на першому рядку:\par
курс, предмет, кількість пропусків;

{\noindent}--- на наступних рядках, починаючи з табуляції, вивести пропуски предмета (по одному на рядок):\par
прізвище, ім’я, кількість пропусків лекцій

{\noindent}у такому сортуванні: прізвище, ім’я.

%Предмети сортувати: курс, предмет.

\bigskip\textit{Обмеження}

Навчальний тиждень може приймати цілі значення від 1 до 19.
День тижня є цілим від 1 до 5. Пара є цілим від 1 до 4.
Курс є цілим від 1 до 4. Аудиторія є додатнім цілим.
Вид занять може приймати значення:
L --- лекція, Пр --- практичне, Seminar --- семінар, ЛАБ --- лабораторне.

Решта полів є непорожніми рядками.
Рядки починатися та закінчуватися білими символами не можуть.

Групи кодуються в межах курсу.
Код групи складається не більш ніж з 4 символів.
Може містити літери алфавіту, десяткові цифри та дефіс.

Предмет є рядком, може мати від 4 до 28 символів включно.
Може містити літери алфавіту, десяткові цифри, апостроф, дефіс, пробіл.
Предмети різних курсів вважати різними навіть за однакових назв.

Прізвище, ім’я та по-батькові (за наявності у варіанті)
можуть мати від 4 до 27, 24, 22 відповідно символів включно кожне.
Можуть містити літери алфавіту, апостроф, пробіл та дефіс.
Студенти однозначно ідентифікуються прізвищем, ім’ям та по-батькові (за наявності у варіанті).
Жоден студент не вчиться одночасно на різних курсах чи в різних групах.

Види занять сортуються так: лекція < лабораторне < семінар < практичне.

%\bigskip%\textit{Для деяких варіантів може знадобитись}


\par
\newpage
{\center Варіант  \No \textbf { 14 }\par}
\bigskip
%type = absense
%variant = 34
Обробляються дані семестрового журналу перевірок відвідування занять студентами деканатом (фіксуються відсутні). 

\underline{Записи основного файлу містять поля}: предмет, вид занять, аудиторія, прізвище, навчальний тиждень, день тижня, пара, ім’я, курс.

\underline{У допоміжному файлі наявні ключі}: сума пар, найбільша довжина назви предмета.

\underline{Знайти} всіх студентів, що пропустили найбільше практичних занять. Вивести по кожному з них інформацію:\par
{\noindent}--- на першому рядку:\par
кількість пропусків практичних, прізвище, ім’я, кількість пропущених пар;

{\noindent}--- на наступних рядках, починаючи з табуляції, вивести пропуски студента (по одному на рядок):\par
предмет, тиждень, кількість пропусків, вид занять

{\noindent}у такому сортуванні: предмет, тиждень, вид занять, кількість пропусків.

%Студентів сортувати: кількість пропусків практичних (за спаданням), прізвище, ім’я.

\bigskip\textit{Обмеження}

Навчальний тиждень може приймати цілі значення від 1 до 15.
День тижня є цілим від 1 до 5. Пара є цілим від 1 до 4.
Курс є цілим від 1 до 4. Аудиторія є додатнім цілим.
Вид занять може приймати значення:
lect --- лекція, P --- практичне, S --- семінар, Lab --- лабораторне.

Решта полів є непорожніми рядками.
Рядки починатися та закінчуватися білими символами не можуть.

Групи кодуються в межах курсу.
Код групи складається не більш ніж з 5 символів.
Може містити літери алфавіту, десяткові цифри та дефіс.

Предмет є рядком, може мати від 3 до 30 символів включно.
Може містити літери алфавіту, десяткові цифри, апостроф, дефіс, пробіл.
Предмети різних курсів вважати різними навіть за однакових назв.

Прізвище, ім’я та по-батькові (за наявності у варіанті)
можуть мати від 3 до 28, 22, 28 відповідно символів включно кожне.
Можуть містити літери алфавіту, апостроф, пробіл та дефіс.
Студенти однозначно ідентифікуються прізвищем, ім’ям та по-батькові (за наявності у варіанті).
Жоден студент не вчиться одночасно на різних курсах чи в різних групах.

Види занять сортуються так: лекція < лабораторне < практичне < семінар.

%\bigskip%\textit{Для деяких варіантів може знадобитись}


\par
\newpage
{\center Варіант  \No \textbf { 15 }\par}
\bigskip
%type = absense
%variant = 40
Обробляються дані семестрового журналу перевірок відвідування занять студентами деканатом (фіксуються відсутні). 

\underline{Записи основного файлу містять поля}: аудиторія, вид занять, курс, ім’я, предмет, день тижня, пара, прізвище, навчальний тиждень.

\underline{У допоміжному файлі наявні ключі}: довжина найдовшої назви предмету, загальна кількість записів.

\underline{Знайти} всі предмети, які пропускалися, але пропусків саме лабораторних занять зафіксовано не було. Вивести по кожному з них інформацію:\par
{\noindent}--- на першому рядку:\par
курс, предмет, кількість пропусків;

{\noindent}--- на наступних рядках, починаючи з табуляції, вивести пропуски предмета (по одному на рядок):\par
тиждень, день, вид занять, кількість пропусків

{\noindent}у такому сортуванні: тиждень, вид занять, день.

%Премети сортувати: курс, кількість пропусків (за спаданням), предмет.

\bigskip\textit{Обмеження}

Навчальний тиждень може приймати цілі значення від 1 до 15.
День тижня є цілим від 1 до 5. Пара є цілим від 1 до 4.
Курс є цілим від 1 до 4. Аудиторія є додатнім цілим.
Вид занять може приймати значення:
Lect --- лекція, Practice --- практичне, Сем. --- семінар, Лаб. --- лабораторне.

Решта полів є непорожніми рядками.
Рядки починатися та закінчуватися білими символами не можуть.

Групи кодуються в межах курсу.
Код групи складається не більш ніж з 6 символів.
Може містити літери алфавіту, десяткові цифри та дефіс.

Предмет є рядком, може мати від 5 до 24 символів включно.
Може містити літери алфавіту, десяткові цифри, апостроф, дефіс, пробіл.
Предмети різних курсів вважати різними навіть за однакових назв.

Прізвище, ім’я та по-батькові (за наявності у варіанті)
можуть мати від 5 до 20, 29, 25 відповідно символів включно кожне.
Можуть містити літери алфавіту, апостроф, пробіл та дефіс.
Студенти однозначно ідентифікуються прізвищем, ім’ям та по-батькові (за наявності у варіанті).
Жоден студент не вчиться одночасно на різних курсах чи в різних групах.

Види занять сортуються так: семінар < лекція < практичне < лабораторне.

%\bigskip%\textit{Для деяких варіантів може знадобитись}


\par
\newpage
{\center Варіант  \No \textbf { 16 }\par}
\bigskip
%type = absense
%variant = 44
Обробляються дані семестрового журналу перевірок відвідування занять студентами деканатом (фіксуються відсутні). 

\underline{Записи основного файлу містять поля}: курс, код групи, день тижня, пара, предмет, вид занять, по-батькові, прізвище, ім’я, аудиторія, навчальний тиждень.

\underline{У допоміжному файлі наявні ключі}: загальна кількість предметів, кількість пропусків, зафіксованих на 4-му курсі.

\underline{Знайти} всі предмети, з яких пропускалися четверті пари. Вивести по кожному з них інформацію:\par
{\noindent}--- на першому рядку:\par
курс, предмет, кількість пропусків, кількість пропусків четвертих пар;

{\noindent}--- на наступних рядках, починаючи з табуляції, вивести пропуски четвертих пар (по одному на рядок):\par
прізвище, ім’я, по-батькові, день, кількість пропусків

{\noindent}у такому сортуванні: день, прізвище, ім’я, по-батькові.

%Предмети сортувати: кількість пропусків, курс, предмет.

\bigskip\textit{Обмеження}

Навчальний тиждень може приймати цілі значення від 1 до 18.
День тижня є цілим від 1 до 5. Пара є цілим від 1 до 4.
Курс є цілим від 1 до 4. Аудиторія є додатнім цілим.
Вид занять може приймати значення:
ЛЕ --- лекція, p --- практичне, sem --- семінар, Lab --- лабораторне.

Решта полів є непорожніми рядками.
Рядки починатися та закінчуватися білими символами не можуть.

Групи кодуються в межах курсу.
Код групи складається не більш ніж з 5 символів.
Може містити літери алфавіту, десяткові цифри та дефіс.

Предмет є рядком, може мати від 6 до 20 символів включно.
Може містити літери алфавіту, десяткові цифри, апостроф, дефіс, пробіл.
Предмети різних курсів вважати різними навіть за однакових назв.

Прізвище, ім’я та по-батькові (за наявності у варіанті)
можуть мати від 6 до 30, 20, 21 відповідно символів включно кожне.
Можуть містити літери алфавіту, апостроф, пробіл та дефіс.
Студенти однозначно ідентифікуються прізвищем, ім’ям та по-батькові (за наявності у варіанті).
Жоден студент не вчиться одночасно на різних курсах чи в різних групах.

Види занять сортуються так: лекція < практичне < лабораторне < семінар.

%\bigskip%\textit{Для деяких варіантів може знадобитись}


\par
\newpage
{\center Варіант  \No \textbf { 17 }\par}
\bigskip
%type = absense
%variant = 39
Обробляються дані семестрового журналу перевірок відвідування занять студентами деканатом (фіксуються відсутні). 

\underline{Записи основного файлу містять поля}: аудиторія, навчальний тиждень, курс, код групи, предмет, вид занять, прізвище, ім’я, день тижня, пара.

\underline{У допоміжному файлі наявні ключі}: загальна кількість записів, FINE.

\underline{Знайти} всі предмети, що пропускалися більше, ніж в середньому. Вивести по кожному з них інформацію:\par
{\noindent}--- на першому рядку:\par
кількість пропусків, курс, предмет, кількість пропусків лабораторних;

{\noindent}--- на наступних рядках, починаючи з табуляції, вивести пропуски предмета (по одному на рядок):\par
прізвище, ім’я, тиждень, вид занять, кількість пропусків

{\noindent}у такому сортуванні: прізвище, ім’я, вид занять.

%Предмети сортувати: кількість пропусків (за спаданням), курс, предмет.

\bigskip\textit{Обмеження}

Навчальний тиждень може приймати цілі значення від 1 до 19.
День тижня є цілим від 1 до 5. Пара є цілим від 1 до 4.
Курс є цілим від 1 до 4. Аудиторія є додатнім цілим.
Вид занять може приймати значення:
лекц. --- лекція, Практ. --- практичне, Сем. --- семінар, лаб. --- лабораторне.

Решта полів є непорожніми рядками.
Рядки починатися та закінчуватися білими символами не можуть.

Групи кодуються в межах курсу.
Код групи складається не більш ніж з 4 символів.
Може містити літери алфавіту, десяткові цифри та дефіс.

Предмет є рядком, може мати від 6 до 33 символів включно.
Може містити літери алфавіту, десяткові цифри, апостроф, дефіс, пробіл.
Предмети різних курсів вважати різними навіть за однакових назв.

Прізвище, ім’я та по-батькові (за наявності у варіанті)
можуть мати від 6 до 21, 21, 26 відповідно символів включно кожне.
Можуть містити літери алфавіту, апостроф, пробіл та дефіс.
Студенти однозначно ідентифікуються прізвищем, ім’ям та по-батькові (за наявності у варіанті).
Жоден студент не вчиться одночасно на різних курсах чи в різних групах.

Види занять сортуються так: практичне < лабораторне < семінар < лекція.

%\bigskip%\textit{Для деяких варіантів може знадобитись}


\par
\newpage
{\center Варіант  \No \textbf { 18 }\par}
\bigskip
%type = absense
%variant = 31
Обробляються дані семестрового журналу перевірок відвідування занять студентами деканатом (фіксуються відсутні). 

\underline{Записи основного файлу містять поля}: предмет, вид занять, ім’я, курс, код групи, день тижня, пара, прізвище, навчальний тиждень, аудиторія.

\underline{У допоміжному файлі наявні ключі}: кількість пропусків на четвертих парах, кількість пропусків лекцій.

\underline{Знайти} всіх студентів, що пропускали заняття виключно на четвертих парах. Вивести по кожному з них інформацію:\par
{\noindent}--- на першому рядку:\par
прізвище, ім’я, кількість пропущених пар;

{\noindent}--- на наступних рядках, починаючи з табуляції, вивести пропуски студента (по одному на рядок):\par
тиждень, день, пара, кількість пропусків

{\noindent}у такому сортуванні: тиждень, день, пара.

%Студентів сортувати: кількість пропусків, прізвище, ім’я.

\bigskip\textit{Обмеження}

Навчальний тиждень може приймати цілі значення від 1 до 20.
День тижня є цілим від 1 до 5. Пара є цілим від 1 до 4.
Курс є цілим від 1 до 4. Аудиторія є додатнім цілим.
Вид занять може приймати значення:
Le --- лекція, pr --- практичне, с. --- семінар, Лаб. --- лабораторне.

Решта полів є непорожніми рядками.
Рядки починатися та закінчуватися білими символами не можуть.

Групи кодуються в межах курсу.
Код групи складається не більш ніж з 6 символів.
Може містити літери алфавіту, десяткові цифри та дефіс.

Предмет є рядком, може мати від 3 до 37 символів включно.
Може містити літери алфавіту, десяткові цифри, апостроф, дефіс, пробіл.
Предмети різних курсів вважати різними навіть за однакових назв.

Прізвище, ім’я та по-батькові (за наявності у варіанті)
можуть мати від 3 до 23, 27, 29 відповідно символів включно кожне.
Можуть містити літери алфавіту, апостроф, пробіл та дефіс.
Студенти однозначно ідентифікуються прізвищем, ім’ям та по-батькові (за наявності у варіанті).
Жоден студент не вчиться одночасно на різних курсах чи в різних групах.

Види занять сортуються так: лабораторне < семінар < практичне < лекція.

%\bigskip%\textit{Для деяких варіантів може знадобитись}


\par
\newpage
{\center Варіант  \No \textbf { 19 }\par}
\bigskip
%type = absense
%variant = 41
Обробляються дані семестрового журналу перевірок відвідування занять студентами деканатом (фіксуються відсутні). 

\underline{Записи основного файлу містять поля}: по-батькові, вид занять, ім’я, предмет, курс, код групи, день тижня, пара, навчальний тиждень, прізвище, аудиторія.

\underline{У допоміжному файлі наявні ключі}: загальна кількість записів, назва предмету, що пропускався найбільше.

\underline{Знайти} всі предмети, що пропускалися тільки на перших трьох навчальних тижнях. Вивести по кожному з них інформацію:\par
{\noindent}--- на першому рядку:\par
курс, предмет, кількість пропусків на першому тижні;

{\noindent}--- на наступних рядках, починаючи з табуляції, вивести пропуски предмета (по одному на рядок):\par
тиждень, день, пара, вид занять, прізвище, ім’я, по-батькові

{\noindent}у такому сортуванні: тиждень, прізвище, ім’я, по-батькові, день, пара, вид занять.

%Предмети сортувати: курс, кількість пропусків(за спаданням), предмет.

\bigskip\textit{Обмеження}

Навчальний тиждень може приймати цілі значення від 1 до 17.
День тижня є цілим від 1 до 5. Пара є цілим від 1 до 4.
Курс є цілим від 1 до 4. Аудиторія є додатнім цілим.
Вид занять може приймати значення:
lect --- лекція, ПР --- практичне, S --- семінар, ЛАБ --- лабораторне.

Решта полів є непорожніми рядками.
Рядки починатися та закінчуватися білими символами не можуть.

Групи кодуються в межах курсу.
Код групи складається не більш ніж з 5 символів.
Може містити літери алфавіту, десяткові цифри та дефіс.

Предмет є рядком, може мати від 4 до 40 символів включно.
Може містити літери алфавіту, десяткові цифри, апостроф, дефіс, пробіл.
Предмети різних курсів вважати різними навіть за однакових назв.

Прізвище, ім’я та по-батькові (за наявності у варіанті)
можуть мати від 4 до 26, 25, 27 відповідно символів включно кожне.
Можуть містити літери алфавіту, апостроф, пробіл та дефіс.
Студенти однозначно ідентифікуються прізвищем, ім’ям та по-батькові (за наявності у варіанті).
Жоден студент не вчиться одночасно на різних курсах чи в різних групах.

Види занять сортуються так: семінар < лекція < практичне < лабораторне.

%\bigskip%\textit{Для деяких варіантів може знадобитись}


\par
\newpage
{\center Варіант  \No \textbf { 20 }\par}
\bigskip
%type = absense
%variant = 42
Обробляються дані семестрового журналу перевірок відвідування занять студентами деканатом (фіксуються відсутні). 

\underline{Записи основного файлу містять поля}: курс, код групи, предмет, вид занять, прізвище, ім’я, навчальний тиждень, день тижня, пара, аудиторія.

\underline{У допоміжному файлі наявні ключі}: кількість різних предметів, кількість записів.

\underline{Знайти} всі предмети, заняття з яких пропускалися тільки в аудиторіях 18 та 42. Вивести по кожному з них інформацію:\par
{\noindent}--- на першому рядку:\par
курс, предмет, кількість пропусків, кількість пропусків лабораторних;

{\noindent}--- на наступних рядках, починаючи з табуляції, вивести пропуски предмета (по одному на рядок):\par
тиждень, день, пара, аудиторія, вид занять, кількість пропусків

{\noindent}у такому сортуванні: тиждень, день, пара, аудиторія, вид занять.

%Предмети сортувати: курс, предмет.

\bigskip\textit{Обмеження}

Навчальний тиждень може приймати цілі значення від 1 до 16.
День тижня є цілим від 1 до 5. Пара є цілим від 1 до 4.
Курс є цілим від 1 до 4. Аудиторія є додатнім цілим.
Вид занять може приймати значення:
Lecture --- лекція, практ. --- практичне, s --- семінар, lab --- лабораторне.

Решта полів є непорожніми рядками.
Рядки починатися та закінчуватися білими символами не можуть.

Групи кодуються в межах курсу.
Код групи складається не більш ніж з 6 символів.
Може містити літери алфавіту, десяткові цифри та дефіс.

Предмет є рядком, може мати від 5 до 29 символів включно.
Може містити літери алфавіту, десяткові цифри, апостроф, дефіс, пробіл.
Предмети різних курсів вважати різними навіть за однакових назв.

Прізвище, ім’я та по-батькові (за наявності у варіанті)
можуть мати від 5 до 22, 23, 20 відповідно символів включно кожне.
Можуть містити літери алфавіту, апостроф, пробіл та дефіс.
Студенти однозначно ідентифікуються прізвищем, ім’ям та по-батькові (за наявності у варіанті).
Жоден студент не вчиться одночасно на різних курсах чи в різних групах.

Види занять сортуються так: практичне < лекція < лабораторне < семінар.

%\bigskip%\textit{Для деяких варіантів може знадобитись}


\par
\newpage
{\center Варіант  \No \textbf { 21 }\par}
\bigskip
%type = absense
%variant = 45
Обробляються дані семестрового журналу перевірок відвідування занять студентами деканатом (фіксуються відсутні). 

\underline{Записи основного файлу містять поля}: прізвище, курс, код групи, ім’я, аудиторія, вид занять, предмет, день тижня, пара, навчальний тиждень.

\underline{У допоміжному файлі наявні ключі}: загальна кількість предметів, кількість пропусків лабораторних занять.

\underline{Знайти} всі предмети, що пропускалися кожного навчального тижня. Вивести по кожному з них інформацію:\par
{\noindent}--- на першому рядку:\par
кількість пропусків, предмет, курс;

{\noindent}--- на наступних рядках, починаючи з табуляції, вивести пропуски предмета (по одному на рядок):\par
тиждень, вид занять, кількість пропусків

{\noindent}у такому сортуванні: тиждень, вид занять.

%Предмети сортувати: кількість пропусків, предмет, курс.

\bigskip\textit{Обмеження}

Навчальний тиждень може приймати цілі значення від 1 до 14.
День тижня є цілим від 1 до 5. Пара є цілим від 1 до 4.
Курс є цілим від 1 до 4. Аудиторія є додатнім цілим.
Вид занять може приймати значення:
Лекц. --- лекція, P --- практичне, сем. --- семінар, ЛАБ --- лабораторне.

Решта полів є непорожніми рядками.
Рядки починатися та закінчуватися білими символами не можуть.

Групи кодуються в межах курсу.
Код групи складається не більш ніж з 4 символів.
Може містити літери алфавіту, десяткові цифри та дефіс.

Предмет є рядком, може мати від 6 до 26 символів включно.
Може містити літери алфавіту, десяткові цифри, апостроф, дефіс, пробіл.
Предмети різних курсів вважати різними навіть за однакових назв.

Прізвище, ім’я та по-батькові (за наявності у варіанті)
можуть мати від 6 до 24, 28, 24 відповідно символів включно кожне.
Можуть містити літери алфавіту, апостроф, пробіл та дефіс.
Студенти однозначно ідентифікуються прізвищем, ім’ям та по-батькові (за наявності у варіанті).
Жоден студент не вчиться одночасно на різних курсах чи в різних групах.

Види занять сортуються так: семінар < практичне < лабораторне < лекція.

%\bigskip%\textit{Для деяких варіантів може знадобитись}


\par
\newpage
{\center Варіант  \No \textbf { 22 }\par}
\bigskip
%type = absense
%variant = 26
Обробляються дані семестрового журналу перевірок відвідування занять студентами деканатом (фіксуються відсутні). 

\underline{Записи основного файлу містять поля}: предмет, навчальний тиждень, вид занять, ім’я, прізвище, курс, код групи, день тижня, пара, по-батькові, аудиторія.

\underline{У допоміжному файлі наявні ключі}: найбільша кількість пропусків одного студента, загальна кількість записів.

\underline{Знайти} всіх студентів, що кожного навчального тижня пропускали хоча б по одній парі. Вивести по кожному з них інформацію:\par
{\noindent}--- на першому рядку:\par
кількість пропущених пар, прізвище, ім’я по-батькові, середня кількість пропущених пар на тиждень (округлена до 1 знаку);

{\noindent}--- на наступних рядках, починаючи з табуляції, вивести пропуски студента (по одному на рядок):\par
тиждень, предмет, кількість пропусків предмета на цьому тижні

{\noindent}у такому сортуванні: тиждень, предмет.

%Студентів сортувати: кількість пропусків(за спаданням), кількість пропусків лекцій, прізвище, ім’я, по-батькові.

\bigskip\textit{Обмеження}

Навчальний тиждень може приймати цілі значення від 1 до 18.
День тижня є цілим від 1 до 5. Пара є цілим від 1 до 4.
Курс є цілим від 1 до 4. Аудиторія є додатнім цілим.
Вид занять може приймати значення:
ЛЕ --- лекція, п --- практичне, Seminar --- семінар, lab --- лабораторне.

Решта полів є непорожніми рядками.
Рядки починатися та закінчуватися білими символами не можуть.

Групи кодуються в межах курсу.
Код групи складається не більш ніж з 5 символів.
Може містити літери алфавіту, десяткові цифри та дефіс.

Предмет є рядком, може мати від 3 до 36 символів включно.
Може містити літери алфавіту, десяткові цифри, апостроф, дефіс, пробіл.
Предмети різних курсів вважати різними навіть за однакових назв.

Прізвище, ім’я та по-батькові (за наявності у варіанті)
можуть мати від 3 до 29, 26, 23 відповідно символів включно кожне.
Можуть містити літери алфавіту, апостроф, пробіл та дефіс.
Студенти однозначно ідентифікуються прізвищем, ім’ям та по-батькові (за наявності у варіанті).
Жоден студент не вчиться одночасно на різних курсах чи в різних групах.

Види занять сортуються так: лекція < практичне < семінар < лабораторне.

%\bigskip%\textit{Для деяких варіантів може знадобитись}


\par
\newpage
{\center Варіант  \No \textbf { 23 }\par}
\bigskip
%type = absense
%variant = 33
Обробляються дані семестрового журналу перевірок відвідування занять студентами деканатом (фіксуються відсутні). 

\underline{Записи основного файлу містять поля}: по-батькові, прізвище, день тижня, пара, вид занять, курс, код групи, аудиторія, навчальний тиждень, ім’я, предмет.

\underline{У допоміжному файлі наявні ключі}: загальна кількість записів, довжина найкоротшого прізвища.

\underline{Знайти} всіх студентів, що пропускали заняття принаймні з двох предметів. Вивести по кожному з них інформацію:\par
{\noindent}--- на першому рядку:\par
прізвище, ім’я, по-батькові, кількість пропущених лекцій;

{\noindent}--- на наступних рядках, починаючи з табуляції, вивести пропуски студента (по одному на рядок):\par
предмет, вид занять, кількість пропусків

{\noindent}у такому сортуванні: предмет, вид занять.

%Студентів сортувати: кількість пропусків лекцій (за спаданням), прізвище, ім’я, по-батькові.

\bigskip\textit{Обмеження}

Навчальний тиждень може приймати цілі значення від 1 до 16.
День тижня є цілим від 1 до 5. Пара є цілим від 1 до 4.
Курс є цілим від 1 до 4. Аудиторія є додатнім цілим.
Вид занять може приймати значення:
лекц. --- лекція, Практ. --- практичне, sem --- семінар, Lab --- лабораторне.

Решта полів є непорожніми рядками.
Рядки починатися та закінчуватися білими символами не можуть.

Групи кодуються в межах курсу.
Код групи складається не більш ніж з 4 символів.
Може містити літери алфавіту, десяткові цифри та дефіс.

Предмет є рядком, може мати від 5 до 23 символів включно.
Може містити літери алфавіту, десяткові цифри, апостроф, дефіс, пробіл.
Предмети різних курсів вважати різними навіть за однакових назв.

Прізвище, ім’я та по-батькові (за наявності у варіанті)
можуть мати від 5 до 30, 30, 20 відповідно символів включно кожне.
Можуть містити літери алфавіту, апостроф, пробіл та дефіс.
Студенти однозначно ідентифікуються прізвищем, ім’ям та по-батькові (за наявності у варіанті).
Жоден студент не вчиться одночасно на різних курсах чи в різних групах.

Види занять сортуються так: лабораторне < практичне < семінар < лекція.

%\bigskip%\textit{Для деяких варіантів може знадобитись}


\par
\newpage
{\center Варіант  \No \textbf { 24 }\par}
\bigskip
%type = absense
%variant = 43
Обробляються дані семестрового журналу перевірок відвідування занять студентами деканатом (фіксуються відсутні). 

\underline{Записи основного файлу містять поля}: день тижня, пара, прізвище, вид занять, предмет, курс, навчальний тиждень, аудиторія, ім’я.

\underline{У допоміжному файлі наявні ключі}: найбільша довжина прізвища, найменша довжина прізвища, середнє значення курсу з округленням до 1 знаку.

\underline{Знайти} всі предмети, з яких пропускалися тільки перші пари. Вивести по кожному з них інформацію:\par
{\noindent}--- на першому рядку:\par
курс, предмет, кількість пропусків;

{\noindent}--- на наступних рядках, починаючи з табуляції, вивести пропуски предмета (по одному на рядок):\par
день, пара, прізвище, ім’я, кількість пропусків

{\noindent}у такому сортуванні: прізвище, ім’я, день, пара.

%Предмети сортувати: кількість пропусків, курс, предмет.

\bigskip\textit{Обмеження}

Навчальний тиждень може приймати цілі значення від 1 до 15.
День тижня є цілим від 1 до 5. Пара є цілим від 1 до 4.
Курс є цілим від 1 до 4. Аудиторія є додатнім цілим.
Вид занять може приймати значення:
Л --- лекція, pr --- практичне, Sem --- семінар, лаб. --- лабораторне.

Решта полів є непорожніми рядками.
Рядки починатися та закінчуватися білими символами не можуть.

Групи кодуються в межах курсу.
Код групи складається не більш ніж з 6 символів.
Може містити літери алфавіту, десяткові цифри та дефіс.

Предмет є рядком, може мати від 4 до 38 символів включно.
Може містити літери алфавіту, десяткові цифри, апостроф, дефіс, пробіл.
Предмети різних курсів вважати різними навіть за однакових назв.

Прізвище, ім’я та по-батькові (за наявності у варіанті)
можуть мати від 4 до 22, 29, 29 відповідно символів включно кожне.
Можуть містити літери алфавіту, апостроф, пробіл та дефіс.
Студенти однозначно ідентифікуються прізвищем, ім’ям та по-батькові (за наявності у варіанті).
Жоден студент не вчиться одночасно на різних курсах чи в різних групах.

Види занять сортуються так: лабораторне < семінар < практичне < лекція.

%\bigskip%\textit{Для деяких варіантів може знадобитись}


\par
\newpage
{\center Варіант  \No \textbf { 25 }\par}
\bigskip
%type = absense
%variant = 50
Обробляються дані семестрового журналу перевірок відвідування занять студентами деканатом (фіксуються відсутні). 

\underline{Записи основного файлу містять поля}: курс, код групи, по-батькові, вид занять, навчальний тиждень, ім’я, предмет, день тижня, пара, прізвище, аудиторія.

\underline{У допоміжному файлі наявні ключі}: загальна кількість пропусків третіх пар, найменша довжина назв предметів, найбільша довжина назв предметів.

\underline{Знайти} всі предмети, які найбільше пропускалися по понеділках. Вивести по кожному з них інформацію:\par
{\noindent}--- на першому рядку:\par
курс, предмет, середня кількість пропусків по днях тижня (округлена до 2 знаків), кількість пропусків практичних;

{\noindent}--- на наступних рядках, починаючи з табуляції, вивести пропуски предмета (по одному на рядок):\par
прізвище, ім’я, по-батькові, тиждень, день, пара, вид занять

{\noindent}у такому сортуванні: день, прізвище, ім’я, по-батькові, тиждень, пара, вид занять.

%Предмети сортувати: середня кількість пропусків по днях тижня (після округлення, за спаданням), кількість пропусків практичних, курс, предмет.

\bigskip\textit{Обмеження}

Навчальний тиждень може приймати цілі значення від 1 до 20.
День тижня є цілим від 1 до 5. Пара є цілим від 1 до 4.
Курс є цілим від 1 до 4. Аудиторія є додатнім цілим.
Вид занять може приймати значення:
Le --- лекція, Practice --- практичне, Sem --- семінар, Лаб. --- лабораторне.

Решта полів є непорожніми рядками.
Рядки починатися та закінчуватися білими символами не можуть.

Групи кодуються в межах курсу.
Код групи складається не більш ніж з 5 символів.
Може містити літери алфавіту, десяткові цифри та дефіс.

Предмет є рядком, може мати від 5 до 22 символів включно.
Може містити літери алфавіту, десяткові цифри, апостроф, дефіс, пробіл.
Предмети різних курсів вважати різними навіть за однакових назв.

Прізвище, ім’я та по-батькові (за наявності у варіанті)
можуть мати від 5 до 23, 22, 24 відповідно символів включно кожне.
Можуть містити літери алфавіту, апостроф, пробіл та дефіс.
Студенти однозначно ідентифікуються прізвищем, ім’ям та по-батькові (за наявності у варіанті).
Жоден студент не вчиться одночасно на різних курсах чи в різних групах.

Види занять сортуються так: лабораторне < семінар < практичне < лекція.

%\bigskip%\textit{Для деяких варіантів може знадобитись}


\par
\newpage
{\center Варіант  \No \textbf { 26 }\par}
\bigskip
%type = absense
%variant = 36
Обробляються дані семестрового журналу перевірок відвідування занять студентами деканатом (фіксуються відсутні). 

\underline{Записи основного файлу містять поля}: \ навчальний\ тиждень,\ по-батькові,\ вид\ занять,\ прізвище, ім’я, аудиторія, предмет, день тижня, пара, курс, код групи.

\underline{У допоміжному файлі наявні ключі}: довжина найкоротшого прізвища, загальна кількість записів.

\underline{Знайти} всіх студентів, що більше інших пропускали заняття по п’ятницях. Вивести по кожному з них інформацію:\par
{\noindent}--- на першому рядку:\par
кількість пропущених пар, кількість пропущених пар по п’ятницях, прізвище, ім’я, по-батькові;

{\noindent}--- на наступних рядках, починаючи з табуляції, вивести пропуски студента саме по п’ятницях (по одному на рядок):\par
тиждень, пара, аудиторія, вид занять, предмет

{\noindent}у такому сортуванні: тиждень, пара, аудиторія, предмет, вид занять.

%Студентів сортувати: кількість пропущених пар, прізвище, ім’я, по-батькові.

\bigskip\textit{Обмеження}

Навчальний тиждень може приймати цілі значення від 1 до 17.
День тижня є цілим від 1 до 5. Пара є цілим від 1 до 4.
Курс є цілим від 1 до 4. Аудиторія є додатнім цілим.
Вид занять може приймати значення:
L --- лекція, Pr --- практичне, Seminar --- семінар, ЛАБ --- лабораторне.

Решта полів є непорожніми рядками.
Рядки починатися та закінчуватися білими символами не можуть.

Групи кодуються в межах курсу.
Код групи складається не більш ніж з 4 символів.
Може містити літери алфавіту, десяткові цифри та дефіс.

Предмет є рядком, може мати від 4 до 33 символів включно.
Може містити літери алфавіту, десяткові цифри, апостроф, дефіс, пробіл.
Предмети різних курсів вважати різними навіть за однакових назв.

Прізвище, ім’я та по-батькові (за наявності у варіанті)
можуть мати від 4 до 26, 26, 28 відповідно символів включно кожне.
Можуть містити літери алфавіту, апостроф, пробіл та дефіс.
Студенти однозначно ідентифікуються прізвищем, ім’ям та по-батькові (за наявності у варіанті).
Жоден студент не вчиться одночасно на різних курсах чи в різних групах.

Види занять сортуються так: семінар < практичне < лабораторне < лекція.

%\bigskip%\textit{Для деяких варіантів може знадобитись}


\par
\newpage
{\center Варіант  \No \textbf { 27 }\par}
\bigskip
%type = absense
%variant = 48
Обробляються дані семестрового журналу перевірок відвідування занять студентами деканатом (фіксуються відсутні). 

\underline{Записи основного файлу містять поля}: прізвище, навчальний тиждень, аудиторія, предмет, вид занять, ім’я, курс, код групи, день тижня, пара.

\underline{У допоміжному файлі наявні ключі}: кількість пропусків перших пар, загальна кількість записів.

\underline{Знайти} всі предмети, з яких було пропущено найбільше семінарів. Вивести по кожному з них інформацію:\par
{\noindent}--- на першому рядку:\par
курс, предмет, кількість пропусків семінарів, кількість пропусків лекцій, кількість пропусків практичних, кількість пропусків лабораторних занять, кількість пропусків;

{\noindent}--- на наступних рядках, починаючи з табуляції, вивести пропуски семінарів з предмету (по одному на рядок):\par
прізвище, ім’я, по-батькові, тиждень, день, аудиторія, пара

{\noindent}у такому сортуванні: прізвище, ім’я, по-батькові, тиждень, день, пара, аудиторія.

%Предмети сортувати: кількість пропусків (за спаданням), курс, предмет.

\bigskip\textit{Обмеження}

Навчальний тиждень може приймати цілі значення від 1 до 14.
День тижня є цілим від 1 до 5. Пара є цілим від 1 до 4.
Курс є цілим від 1 до 4. Аудиторія є додатнім цілим.
Вид занять може приймати значення:
Lect --- лекція, Пр --- практичне, сем. --- семінар, Lab --- лабораторне.

Решта полів є непорожніми рядками.
Рядки починатися та закінчуватися білими символами не можуть.

Групи кодуються в межах курсу.
Код групи складається не більш ніж з 6 символів.
Може містити літери алфавіту, десяткові цифри та дефіс.

Предмет є рядком, може мати від 3 до 29 символів включно.
Може містити літери алфавіту, десяткові цифри, апостроф, дефіс, пробіл.
Предмети різних курсів вважати різними навіть за однакових назв.

Прізвище, ім’я та по-батькові (за наявності у варіанті)
можуть мати від 3 до 25, 25, 22 відповідно символів включно кожне.
Можуть містити літери алфавіту, апостроф, пробіл та дефіс.
Студенти однозначно ідентифікуються прізвищем, ім’ям та по-батькові (за наявності у варіанті).
Жоден студент не вчиться одночасно на різних курсах чи в різних групах.

Види занять сортуються так: лабораторне < семінар < практичне < лекція.

%\bigskip%\textit{Для деяких варіантів може знадобитись}


\par
\newpage
{\center Варіант  \No \textbf { 28 }\par}
\bigskip
%type = absense
%variant = 23
Обробляються дані семестрового журналу перевірок відвідування занять студентами деканатом (фіксуються відсутні). 

\underline{Записи основного файлу містять поля}: вид занять, аудиторія, навчальний тиждень, день тижня, пара, ім’я, предмет, прізвище, курс.

\underline{У допоміжному файлі наявні ключі}: загальна кількість записів, загальна кількість пропусків семінарів.

\underline{Знайти} всіх студентів, що пропустили більше 10 пар. Вивести по кожному з них інформацію:\par
{\noindent}--- на першому рядку:\par
прізвище, ім’я, кількість пропущених пар, кількість пропущених практичних;

{\noindent}--- на наступних рядках, починаючи з табуляції, вивести пропуски студента (по одному на рядок):\par
предмет, вид занять, тиждень, день, пара, аудиторія

{\noindent}у такому сортуванні: предмет, вид занять, тиждень, день, пара, аудиторія.

%Студентів сортувати: кількість пропусків практичних (за спаданням), кількість пропусків, прізвище, ім’я.

\bigskip\textit{Обмеження}

Навчальний тиждень може приймати цілі значення від 1 до 19.
День тижня є цілим від 1 до 5. Пара є цілим від 1 до 4.
Курс є цілим від 1 до 4. Аудиторія є додатнім цілим.
Вид занять може приймати значення:
Лекц. --- лекція, ПР --- практичне, s --- семінар, Лаб. --- лабораторне.

Решта полів є непорожніми рядками.
Рядки починатися та закінчуватися білими символами не можуть.

Групи кодуються в межах курсу.
Код групи складається не більш ніж з 5 символів.
Може містити літери алфавіту, десяткові цифри та дефіс.

Предмет є рядком, може мати від 6 до 40 символів включно.
Може містити літери алфавіту, десяткові цифри, апостроф, дефіс, пробіл.
Предмети різних курсів вважати різними навіть за однакових назв.

Прізвище, ім’я та по-батькові (за наявності у варіанті)
можуть мати від 6 до 28, 27, 23 відповідно символів включно кожне.
Можуть містити літери алфавіту, апостроф, пробіл та дефіс.
Студенти однозначно ідентифікуються прізвищем, ім’ям та по-батькові (за наявності у варіанті).
Жоден студент не вчиться одночасно на різних курсах чи в різних групах.

Види занять сортуються так: лабораторне < семінар < практичне < лекція.

%\bigskip%\textit{Для деяких варіантів може знадобитись}


\par
\newpage
{\center Варіант  \No \textbf { 29 }\par}
\bigskip
%type = absense
%variant = 29
Обробляються дані семестрового журналу перевірок відвідування занять студентами деканатом (фіксуються відсутні). 

\underline{Записи основного файлу містять поля}: предмет, прізвище, день тижня, пара, аудиторія, вид занять, навчальний тиждень, курс, код групи, ім’я.

\underline{У допоміжному файлі наявні ключі}: найбільший номер аудиторії, кількість записів у файлі.

\underline{Знайти} всіх студентів, що хоча б з одного предмета пропустили 15 лекцій. Вивести по кожному з них інформацію:\par
{\noindent}--- на першому рядку:\par
кількість пропущених пар, кількість пропущених лекцій прізвище, ім’я;

{\noindent}--- на наступних рядках, починаючи з табуляції, вивести пропуски студента (по одному на рядок):\par
предмет, вид занять, кількість пропусків

{\noindent}у такому сортуванні: предмет, вид занять.

%Студентів сортувати: кількість пропусків, прізвище, ім’я.

\bigskip\textit{Обмеження}

Навчальний тиждень може приймати цілі значення від 1 до 19.
День тижня є цілим від 1 до 5. Пара є цілим від 1 до 4.
Курс є цілим від 1 до 4. Аудиторія є додатнім цілим.
Вид занять може приймати значення:
L --- лекція, практ. --- практичне, S --- семінар, лаб. --- лабораторне.

Решта полів є непорожніми рядками.
Рядки починатися та закінчуватися білими символами не можуть.

Групи кодуються в межах курсу.
Код групи складається не більш ніж з 6 символів.
Може містити літери алфавіту, десяткові цифри та дефіс.

Предмет є рядком, може мати від 3 до 34 символів включно.
Може містити літери алфавіту, десяткові цифри, апостроф, дефіс, пробіл.
Предмети різних курсів вважати різними навіть за однакових назв.

Прізвище, ім’я та по-батькові (за наявності у варіанті)
можуть мати від 3 до 29, 20, 25 відповідно символів включно кожне.
Можуть містити літери алфавіту, апостроф, пробіл та дефіс.
Студенти однозначно ідентифікуються прізвищем, ім’ям та по-батькові (за наявності у варіанті).
Жоден студент не вчиться одночасно на різних курсах чи в різних групах.

Види занять сортуються так: лекція < семінар < лабораторне < практичне.

%\bigskip%\textit{Для деяких варіантів може знадобитись}


\par
\newpage
{\center Варіант  \No \textbf { 30 }\par}
\bigskip
%type = absense
%variant = 37
Обробляються дані семестрового журналу перевірок відвідування занять студентами деканатом (фіксуються відсутні). 

\underline{Записи основного файлу містять поля}: курс, код групи, вид занять, навчальний тиждень, предмет, день тижня, пара, аудиторія, прізвище, ім’я.

\underline{У допоміжному файлі наявні ключі}: кількість зафіксованих пропусків, кількість пропусків лекцій.

\underline{Знайти} всіх студентів, що по вівторках пропускали заняття більше ніж у середньому. Вивести по кожному з них інформацію:\par
{\noindent}--- на першому рядку:\par
прізвище, ім’я, середня кількість пропусків (по днях тижня, округлена до 1 знаку), кількість пропусків по вівторках;

{\noindent}--- на наступних рядках, починаючи з табуляції, вивести пропуски студента (по одному на рядок):\par
предмет, вид занять, тиждень, день, пара

{\noindent}у такому сортуванні: день, тиждень, пара, предмет, вид занять.

%Студентів сортувати: середня кількість пропусків (округлена), кількість пропусків по вівторках, прізвище, ім’я.

\bigskip\textit{Обмеження}

Навчальний тиждень може приймати цілі значення від 1 до 14.
День тижня є цілим від 1 до 5. Пара є цілим від 1 до 4.
Курс є цілим від 1 до 4. Аудиторія є додатнім цілим.
Вид занять може приймати значення:
ЛЕ --- лекція, p --- практичне, Сем. --- семінар, lab --- лабораторне.

Решта полів є непорожніми рядками.
Рядки починатися та закінчуватися білими символами не можуть.

Групи кодуються в межах курсу.
Код групи складається не більш ніж з 4 символів.
Може містити літери алфавіту, десяткові цифри та дефіс.

Предмет є рядком, може мати від 6 до 28 символів включно.
Може містити літери алфавіту, десяткові цифри, апостроф, дефіс, пробіл.
Предмети різних курсів вважати різними навіть за однакових назв.

Прізвище, ім’я та по-батькові (за наявності у варіанті)
можуть мати від 6 до 24, 24, 21 відповідно символів включно кожне.
Можуть містити літери алфавіту, апостроф, пробіл та дефіс.
Студенти однозначно ідентифікуються прізвищем, ім’ям та по-батькові (за наявності у варіанті).
Жоден студент не вчиться одночасно на різних курсах чи в різних групах.

Види занять сортуються так: семінар < лекція < практичне < лабораторне.

%\bigskip%\textit{Для деяких варіантів може знадобитись}


\par
\newpage
{\center Варіант  \No \textbf { 31 }\par}
\bigskip
%type = absense
%variant = 47
Обробляються дані семестрового журналу перевірок відвідування занять студентами деканатом (фіксуються відсутні). 

\underline{Записи основного файлу містять поля}: по-батькові, аудиторія, навчальний тиждень, день тижня, пара, предмет, ім’я, вид занять, курс, код групи, прізвище.

\underline{У допоміжному файлі наявні ключі}: середня кількість пропусків предметів, округлена до 2 знаків, SMILE.

\underline{Знайти} всі предмети, більша половина пропусків з яких – пропуски лекцій. Вивести по кожному з них інформацію:\par
{\noindent}--- на першому рядку:\par
курс, предмет, кількість пропусків лекцій, кількість пропусків;

{\noindent}--- на наступних рядках, починаючи з табуляції, вивести пропуски предмета (по одному на рядок):\par
прізвище, ім’я, по-батькові, тиждень, вид занять, кількість пропусків

{\noindent}у такому сортуванні: вид занять, прізвище, ім’я, по-батькові, тиждень.

%Предмети сортувати: кількість пропусків лекцій, курс, предмет.

\bigskip\textit{Обмеження}

Навчальний тиждень може приймати цілі значення від 1 до 20.
День тижня є цілим від 1 до 5. Пара є цілим від 1 до 4.
Курс є цілим від 1 до 4. Аудиторія є додатнім цілим.
Вид занять може приймати значення:
лекц. --- лекція, п --- практичне, sem --- семінар, lab --- лабораторне.

Решта полів є непорожніми рядками.
Рядки починатися та закінчуватися білими символами не можуть.

Групи кодуються в межах курсу.
Код групи складається не більш ніж з 4 символів.
Може містити літери алфавіту, десяткові цифри та дефіс.

Предмет є рядком, може мати від 5 до 24 символів включно.
Може містити літери алфавіту, десяткові цифри, апостроф, дефіс, пробіл.
Предмети різних курсів вважати різними навіть за однакових назв.

Прізвище, ім’я та по-батькові (за наявності у варіанті)
можуть мати від 5 до 20, 21, 27 відповідно символів включно кожне.
Можуть містити літери алфавіту, апостроф, пробіл та дефіс.
Студенти однозначно ідентифікуються прізвищем, ім’ям та по-батькові (за наявності у варіанті).
Жоден студент не вчиться одночасно на різних курсах чи в різних групах.

Види занять сортуються так: семінар < лабораторне < практичне < лекція.

%\bigskip%\textit{Для деяких варіантів може знадобитись}


\par
\newpage
\end{document}
